\begin{tikzpicture}[
    x=1cm,y=1cm,
    >=Latex,
    font=\small,
    titulo/.style={
        draw,
        rounded corners=2pt,
        fill=gray!15,
        font=\bfseries,
        minimum width=3.9cm,
        minimum height=0.8cm,
        align=center
    },
    mineralL/.style={
        draw,
        rounded corners=2pt,
        fill=green!18,
        minimum width=3.6cm,
        minimum height=0.8cm,
        align=center
    },
    mineralR/.style={
        draw,
        rounded corners=2pt,
        fill=blue!12,
        text width=3.25cm,
        minimum height=0.95cm,
        align=center
    },
    tardio/.style={
        draw,
        rounded corners=2pt,
        fill=orange!12,
        minimum width=2.35cm,
        minimum height=0.8cm,
        align=center
    },
    notaL/.style={
        align=center,
        text width=3.0cm,
        font=\footnotesize
    },
    notaR/.style={
        align=center,
        text width=3.3cm,
        font=\footnotesize
    }
]

% -------------------------------------------------
% Título geral
% -------------------------------------------------
\node[font=\bfseries\large] at (0.1,6.0) {Série de Reação de Bowen};

% -------------------------------------------------
% Eixo de temperatura
% -------------------------------------------------
\draw[very thick,-{Latex[length=3mm]}] (-5.4,5.5) -- (-5.4,-5.4);
\node[left] at (-5.55,5.1) {alta};
\node[left] at (-5.55,-4.9) {baixa};
\node[below,align=center] at (-5.4,-5.45) {Temperatura\\decrescente};

% -------------------------------------------------
% Cabeçalhos das colunas
% -------------------------------------------------
\node[titulo] (td) at (-1.9,5.1) {Série descontínua};
\node[titulo] (tc) at ( 2.1,5.1) {Série contínua};

% -------------------------------------------------
% Série descontínua
% -------------------------------------------------
\node[mineralL] (ol) at (-1.9,4.0) {Olivina};
\node[mineralL] (px) at (-1.9,2.8) {Piroxênio};
\node[mineralL] (am) at (-1.9,1.6) {Anfibólio};
\node[mineralL] (bt) at (-1.9,0.4) {Biotita};

\draw[very thick,-{Latex[length=3mm]}] (ol.south) -- (px.north);
\draw[very thick,-{Latex[length=3mm]}] (px.south) -- (am.north);
\draw[very thick,-{Latex[length=3mm]}] (am.south) -- (bt.north);

% -------------------------------------------------
% Série contínua
% -------------------------------------------------
\node[mineralR] (pl1) at (2.1,4.0) {Plagioclásio\\cálcico};
\node[mineralR] (pl2) at (2.1,2.8) {Plagioclásio\\cálcico--sódico};
\node[mineralR] (pl3) at (2.1,1.6) {Plagioclásio\\sódico--cálcico};
\node[mineralR] (pl4) at (2.1,0.4) {Plagioclásio\\sódico};

\draw[very thick,-{Latex[length=3mm]}] (pl1.south) -- (pl2.north);
\draw[very thick,-{Latex[length=3mm]}] (pl2.south) -- (pl3.north);
\draw[very thick,-{Latex[length=3mm]}] (pl3.south) -- (pl4.north);

% -------------------------------------------------
% Notas explicativas (deslocadas para não cruzar setas)
% -------------------------------------------------
\node[notaL] at (-3.5,-0.9)
{cada mineral reage com o magma e é \\substituído por \\outro};

\node[notaR] at (4.1,-0.9)
{a composição muda de rica em Ca\\para rica em Na};

% -------------------------------------------------
% Caixa central: minerais tardios
% -------------------------------------------------
\node[titulo, minimum width=5.5cm] (tf) at (0.1,-1.8) {Minerais tardios};

% -------------------------------------------------
% Conexões até a caixa central (desviadas do texto)
% -------------------------------------------------
\coordinate (ljoin) at (-0.95,-1.15);
\coordinate (rjoin) at ( 1.15,-1.15);

\draw[very thick]
    (bt.south) -- ++(0,-0.15) coordinate (btA)
    -- ++(0.95,0) coordinate (btB)
    -- (btB |- ljoin);
\draw[very thick,-{Latex[length=3mm]}]
    (btB |- ljoin) -- (tf.north west);

\draw[very thick]
    (pl4.south) -- ++(0,-0.15) coordinate (plA)
    -- ++(-0.95,0) coordinate (plB)
    -- (plB |- rjoin);
\draw[very thick,-{Latex[length=3mm]}]
    (plB |- rjoin) -- (tf.north east);

% -------------------------------------------------
% Minerais tardios finais
% -------------------------------------------------
\node[tardio] (kf) at (-2.35,-3.45) {K-feldspato};
\node[tardio] (ms) at ( 0.00,-3.45) {Muscovita};
\node[tardio] (qt) at ( 2.35,-3.45) {Quartzo};

% pontos de saída da caixa central
\coordinate (tfL) at ($(tf.south west)!0.18!(tf.south east)$);
\coordinate (tfC) at ($(tf.south west)!0.50!(tf.south east)$);
\coordinate (tfR) at ($(tf.south west)!0.82!(tf.south east)$);

\draw[very thick,-{Latex[length=3mm]}] (tfL) -- (kf.north);
\draw[very thick,-{Latex[length=3mm]}] (tfC) -- (ms.north);
\draw[very thick,-{Latex[length=3mm]}] (tfR) -- (qt.north);

% -------------------------------------------------
% Legendas laterais afastadas do desenho
% -------------------------------------------------
\node[align=left, text width=2.6cm] at (6.0,4.0)
{minerais de\\alta temperatura};

\node[align=left, text width=2.6cm] at (5.0,-3.45)
{minerais de\\baixa temperatura};

\end{tikzpicture}